\begin{proof}
(i) The $w_{j}$ anticommute and are null (Fact \claimref{field-theory.imported-tower-limit}), so for $x = \sum_{j}c_{j}w_{j}$, $xx = \sum_{j}c_{j}^{2}Q(w_{j}) + \sum_{j<k}c_{j}c_{k}\{w_{j}, w_{k}\} = 0$, and $w_{j}w_{k} = -w_{k}w_{j}$: the products of two annihilators are the $\binom{n}{2}$ elements $w_{j}w_{k}$, $j < k$, linearly independent because they are distinct monomials of the exterior algebra of the null space. The stored block of Proposition \claimref{field-theory.composite-scalars-mean-field-level}(1) has mass term $\mu w_{a}w_{c}$ plus adjoint, the pair term of the two modes of one array (the check recovers its $\mathrm{ad}_{H}^{2} = K^{2} + \mu^{2}$). A bilinear, self-adjoint term that changes the number of quanta by two is a combination of pair terms and their adjoints; between two arrays the pair terms that respect the momentum pairing of the modes ($a$ with a $c$) are $w_{0}w_{1}$, $w_{2}w_{3}$ (self) and $w_{0}w_{3}$, $w_{2}w_{1}$ (cross).

(ii) From $\{v_{j}, w_{k}\} = \delta_{jk}$, $[n_{j}, w_{k}] = -\delta_{jk}w_{k}$ and $[n_{j}, v_{k}] = \delta_{jk}v_{k}$, so a charge $\sum_{j}q_{j}n_{j}$ acts on $w_{j}w_{k}$ by $-(q_{j} + q_{k})$. This gives the table: $N$: $-2$ on all; $N_{a}$ and $N_{c}$: $-1$ on $w_{0}w_{3}$, $w_{2}w_{1}$, $w_{2}w_{3}$ (each has one $a$ mode and one $c$ mode); $L$: $0$ on $w_{0}w_{3}$ and $w_{2}w_{1}$, $+2$ on $w_{2}w_{3}$. If $[N_{a}, X] = -X$ and $p$ is the spectral projection of $N_{a}$ at a value $\lambda$, then $N_{a}Xp = (\lambda - 1)Xp$, so $pXp = 0$; the same for the adjoint with $+1$. The element $H_{4}$ of Theorem \claimref{synthesis.tower-pairing-gap-at-fixed-charges-is-twice-the-quasiparticle-energy} conserves $N_{a}$ and $N_{c}$ and its clauses (i), (ii) are statements on their joint sectors, so a bilinear mass term has no part in them. $\Gamma$ anticommutes with every $v_{j}$, $w_{j}$ and so commutes with every product of two. The commutation statements for $H$ follow from the table, $D$ and $S$ being linearly independent of each other and of their adjoints. In Proposition \claimref{standard-model.majorana-element}(2), $T$ commutes with the gauge algebra and $[B - L, T] = -2T$; here $D$ commutes with $L$ and $S$ does not.

(iii) $\theta$ is induced by an orthogonal map of the generating space (a product of two reflections), so it is an automorphism, and it commutes with reversion and with $f_{j} \mapsto -f_{j}$, so with the adjoint. $\theta(w_{j}) = \tfrac12(e_{j} - f_{j}) = v_{j}$ for $j = 1, 3$, so $\theta(n_{j}) = w_{j}v_{j} = 1 - n_{j}$ there, $\theta(N_{a} - N_{c}) = n_{0} + n_{2} - (1 - n_{1}) - (1 - n_{3}) = N - 2$, and $\theta(S) = w_{2}v_{3}$, $\theta(D) = w_{0}v_{3} + w_{2}v_{1}$, each a product of one annihilator and one creator. $H$ commutes with $N_{a} - N_{c}$ by (ii), so $\theta(H)$ commutes with $N$. An automorphism preserves the minimal polynomial of $H$, so the spectra agree in every faithful module.

(iv) Write the bilinear part as $\tfrac12\sum_{j,k}\Omega_{jk}w_{j}w_{k}$ plus its adjoint $\tfrac12\sum_{j,k}\Omega_{jk}v_{k}v_{j}$, with $\Omega$ the real antisymmetric matrix with $\Omega_{03} = \Omega_{21} = m$, $\Omega_{23} = M$ and the other entries above the diagonal zero. From $[w_{j}w_{k}, v_{l}] = \delta_{kl}w_{j} - \delta_{jl}w_{k}$ and its adjoint, $\mathrm{ad}_{H}v_{l} = \sum_{j}\Omega_{jl}w_{j}$ and $\mathrm{ad}_{H}w_{l} = -\sum_{k}\Omega_{kl}v_{k}$ at $K = 0$, so $\mathrm{ad}_{H}^{2}$ acts on the span of the $v_{l}$ and on the span of the $w_{l}$ by $-\Omega^{2}$. The characteristic polynomial of a real antisymmetric four-by-four matrix is $t^{4} + (\sum_{j<k}\Omega_{jk}^{2})t^{2} + \mathrm{Pf}(\Omega)^{2}$, here $t^{4} + (2m^{2} + M^{2})t^{2} + m^{4}$, since $\mathrm{Pf}(\Omega) = \Omega_{01}\Omega_{23} - \Omega_{02}\Omega_{13} + \Omega_{03}\Omega_{12} = -m^{2}$. By the Cayley-Hamilton identity, $\mathrm{ad}_{H}^{4} - (M^{2} + 2m^{2})\mathrm{ad}_{H}^{2} + m^{4} = 0$ on the mode operators, for all real $m$, $M$; the check confirms it, signs included, at five pairs. The roots in $y = \mathrm{ad}_{H}^{2}$ are the squares of the roots $x_{\pm}$ of $x^{2} - Mx - m^{2}$ (sum of squares $M^{2} + 2m^{2}$, product of squares $m^{4}$), and the eigenvalues of $\mathrm{ad}_{H}$ on mode operators are plus and minus the mode energies.

For $K \neq 0$: $[N, v_{l}] = v_{l}$ and $[N, w_{l}] = -w_{l}$, so in the splitting of the mode operators into the span of the $v_{l}$ and the span of the $w_{l}$, $\mathrm{ad}_{H}$ has diagonal blocks $K$, $-K$ and off-diagonal blocks $-\Omega$, $\Omega$; in the square the cross terms $-K\Omega + \Omega K$ cancel and $\mathrm{ad}_{H}^{2} = K^{2} - \Omega^{2}$ on both spans. Substituting $-\Omega^{2} = \mathrm{ad}_{H}^{2} - K^{2}$ in the Cayley-Hamilton identity gives the stated relation. \emph{Charges conserved by $D$ (clause (ii)).} $\sum_{j}q_{j}n_{j}$ acts on $w_{0}w_{3}$ by $-(q_{0} + q_{3})$ and on $w_{2}w_{1}$ by $-(q_{1} + q_{2})$, so it commutes with $D$ exactly when $q = (\alpha, \beta, -\beta, -\alpha)$, and then acts on $S = w_{2}w_{3}$ by $-(q_{2} + q_{3}) = \alpha + \beta$. \emph{Rank (clause (iii)).} At $m = 0$, $\Omega$ has rank two and $\mathrm{ad}_{H}^{2}$ vanishes on the modes of the first array; at $M = 0$, $-\Omega^{2} = m^{2}$ on all. The stored block: $w_{0}w_{1}$ has $L$ charge $-(1 + 1)$ and $N_{a} - N_{c}$ charge $-(1 - 1)$.

\emph{Reading.} Theorem \claimref{synthesis.light-ray-recurrence-test-verdict-on-generations-is-set-by-the-scalar-content} states its tests through the count $c$ and the grading $P$, and its clause (i) gives one complex array per Weyl component for every mass matrix; (ii) gives $[\Gamma, S] = 0$.
\end{proof}
